\documentclass[10pt,twocolumn]{article}
% --- Packages (imports) ---
\usepackage{amsmath}
\usepackage{graphicx}
\usepackage{geometry}
\usepackage{titlesec}
\usepackage{hyperref}
\usepackage{caption}
\usepackage{booktabs}
\usepackage{array}
\usepackage{cite}
\usepackage{subcaption}
\usepackage{balance}
\usepackage{float}  % Add this for better figure positioning
\usepackage{makecell}
\usepackage[table]{xcolor}
\usepackage{colortbl}
\usepackage{orcidlink}

\definecolor{lightgray}{gray}{0.9}
\definecolor{lightblue}{rgb}{0.85,0.92,1}
\definecolor{lightgreen}{rgb}{0.88,1,0.88}

% Geometry - similar to conference format
\geometry{
  papersize={8.5in,11in},
  top=0.75in,
  bottom=1in,
  left=0.5in,
  right=0.5in,
  columnsep=0.25in
}

% Section formatting - center aligned main headings
\titleformat{\section}{\normalsize\bfseries\centering}{\thesection.}{0.5em}{\MakeUppercase}
\titleformat{\subsection}{\normalsize\bfseries}{\thesubsection}{0.5em}{}

% Hyperref setup - remove boxes around references
\hypersetup{
  colorlinks=false,
  hidelinks,
  pdfborder={0 0 0},
  pdfauthor={Avinash Bhosale, Shaarav Bhosale},
  pdftitle={Key Identification Using Chroma Analysis Of Harmonic Overlays}
}

% Caption formatting to match conference style
\captionsetup{font=small,labelfont=bf}

\newcommand{\safeincludegraphics}[2][]{%
  \includegraphics[#1]{\detokenize{#2}}%
}

% Paragraph formatting - first paragraph no indent, subsequent paragraphs indented
\setlength{\parindent}{1em}  % Standard paragraph indent
\setlength{\parskip}{0pt}    % No space between paragraphs

% Custom command to remove indent for first paragraph after subsection
\makeatletter
\let\@afterindenttrue\@afterindentfalse
\@afterindentfalse
\makeatother

% Title formatting to span full width and be centered at top
\makeatletter
\renewcommand{\maketitle}{
  \twocolumn[
    \begin{center}
      {\large\bfseries\@title\par}
      \vspace{6pt}
      {\normalsize\@author\par}
      \vspace{12pt}
    \end{center}
  ]
}
\makeatother

% Add some microtype for better spacing
\usepackage{microtype}

\title{KEY IDENTIFICATION USING CHROMA ANALYSIS OF HARMONIC OVERLAYS}
\author{Avinash Bhosale\, \orcidlink{0009-0002-4276-8109} Shaarav Bhosale \orcidlink{0009-0009-8207-2887}}

\begin{document}
\maketitle

\begin{abstract}
Determining the key of a musical piece is a fundamental task in Music Information Retrieval (MIR).
Using chroma-based analysis, we want to propose a method which imitates musicians’ key-finding process.
As musicians often hum or play different notes to find the one that resonates best with the track,
we simulate this by generating synthetic harmonic overlays for each note and measuring their tonal coherence with the original audio.
This method creates hybrid signals where each synthetic harmonic is added to the original music and chroma features are extracted for all twelve hybrid signals.
Then using cosine similarity tries to evaluate which hybrid signal’s chroma features best match the original.
This strategy identifies the key by ranking the cosine similarity of each candidate.
However, this method has not been tested for compositions that modulate between multiple keys or exhibit polytonality.
\end{abstract}

\section{Introduction}\label{sec:introduction}
Key detection is a central problem in the field of Music Information Retrieval (MIR), serving as a foundation for tasks such as harmonic analysis, genre classification, and recommendation systems.
Identifying the tonal center of a musical piece allows computational models to capture structural relationships that musicians intuitively perceive, thereby bridging the gap between human listening practices and algorithmic analysis.
Traditional approaches to key detection often rely on statistical distributions of pitch classes or template matching against predefined tonal profiles.
While effective in many cases, these methods can struggle to emulate the dynamic and exploratory strategies employed by musicians themselves.

Musicians frequently determine the key of a piece by actively engaging with the sound—humming, singing, or playing notes until one resonates as the tonal center.
Inspired by this process, we propose a chroma-based method that simulates such human intuition.
Our approach generates synthetic harmonic overlays for each candidate pitch class and combines them with the original audio to form hybrid signals.
By extracting chroma features from these hybrid signals and evaluating their tonal coherence with the original track through cosine similarity, the system identifies the most likely key.
This ranking-based strategy offers a novel perspective on key detection, emphasizing tonal reinforcement rather than static pitch distributions.


\subsection{Challenges Posed by the Perfect Fifth}\label{sec:challenges-posed-by-the-perfect-fifth}

Every basic chord (triad) is built from three notes: the root, the third, and the fifth.
The third determines the chord’s quality (major or minor), while the fifth --- most often a perfect fifth above the root --- provides stability and pleasing to the ear.
Because of this, the perfect fifth is almost universally present in chords, reinforcing its perceptual strength in tonal music.

Even in Indian classical music, the drone sound which exists throughout the performance also has the tonic (Sa) and the perfect fifth (Pa).

Mathematically, a \textbf{major triad} built on a fundamental frequency $f_0$ can be expressed as:
\begin{equation}
\text{Major Triad} = \left\{ f_0,\; f_0 \cdot 2^{\tfrac{4}{12}},\; f_0 \cdot 2^{\tfrac{7}{12}} \right\} \label{eq:major_triad}
\end{equation}

Similarly, a \textbf{minor triad} is given by:
\begin{equation}
\text{Minor Triad} = \left\{ f_0,\; f_0 \cdot 2^{\tfrac{3}{12}},\; f_0 \cdot 2^{\tfrac{7}{12}} \right\} \label{eq:minor_triad}
\end{equation}

Here:
\begin{align*}
f_0 & : \text{root (tonic) frequency} \\
f_0 \cdot 2^{\tfrac{4}{12}} & : \text{major third (4 semitones above root)} \\
f_0 \cdot 2^{\tfrac{3}{12}} & : \text{minor third (3 semitones above root)} \\
f_0 \cdot 2^{\tfrac{7}{12}} & : \text{perfect fifth (7 semitones above root)}
\end{align*}

However, this very strength introduces complications in computational key estimation.
Because the perfect fifth shares strong harmonic energy with the tonic, chroma-based analyses often register elevated similarity between the fifth and the tonic.
In practice, this can lead to \textbf{misclassification}, where the algorithm mistakenly identifies the fifth as the tonal center rather than the intended tonic.
Our initial experiments using tonic-only harmonic overlays highlighted this issue, with accuracy plateauing around 60--70\%.
Chroma CQT plots revealed recurring tonal interference patterns, showing that the fifth note occasionally dominated the ranking process.

This phenomenon underscores the dual role of the perfect fifth: while musically essential, it can obscure algorithmic detection of the true key.


\section{Related Work}\label{sec:related-work}

\subsection{Cognitive Foundations of Tonal Hierarchies}\label{subsec:cognitive-foundations-of-tonal-hierarchies}
Krumhansl's influential work ``Cognitive Foundations of Musical Pitch'' \cite{krumhansl1990cognitive} established the theoretical foundation for computational key identification by demonstrating that listeners internalize abstract tonal hierarchies—systematic orderings of pitch stability reflecting cultural musical knowledge. Through the probe tone method, Krumhansl showed that musical contexts establish hierarchical orderings where certain tones (tonic, dominant, mediant) achieve greater perceptual stability than others.

Crucially, Krumhansl demonstrated that these tonal hierarchies correlate strongly with statistical distributions of tones in musical compositions, providing a mechanism through which abstract pitch relationships are learned from exposure to musical styles. This cognitive framework directly supports our approach: by generating synthetic harmonic overlays, we amplify existing statistical patterns in the audio, leveraging the same perceptual principles that govern human tonal perception.

Krumhansl's cross-cultural studies, including analysis of North Indian ragas, revealed that while the specific hierarchical orderings differ across musical cultures, the general principle of tonal centrality remains universal.

The Krumhansl-Schmuckler algorithm, derived from these cognitive studies, compares pitch class distributions with theoretical key profiles. However, this template-matching approach often fails with perfect fifth ambiguity—the exact problem our method addresses through dual harmonic enhancement.

\subsection{Chroma-Based Harmonic Analysis}\label{subsec:chroma-based-harmonic-analysis}
Chroma features, representing the 12 pitch classes independent of octave, have become standard for harmonic analysis\cite{muller2011chroma}.
The Constant-Q Transform (CQT) provides superior frequency resolution for musical applications compared to traditional FFT-based methods, making it ideal for precise pitch class identification.

Müller and Ewert's comprehensive chroma toolbox demonstrates the critical importance of feature design choices in musical analysis.
Their work reveals that different chroma variants (CP, CLP, CENS, CRP) exhibit vastly different behaviors depending on the analysis task.
Logarithmic compression, temporal smoothing, and normalization strategies all significantly impact performance, emphasizing the need for culturally adaptive feature processing in cross-cultural musical analysis.

\subsection{Cultural Adaptations in MIR}\label{subsec:cultural-adaptations-in-mir}
Limited work exists on adapting key identification algorithms for non-Western musical traditions.
Srinivasamurthy et al. \cite{srinivasamurthy2021saraga}created datasets for Indian classical music analysis, while Senthil Raja~\cite{raja2023sequential} focused specifically on raga identification using sequential pitch distributions.
However, Bollywood music's unique position as a fusion genre requires novel approaches that can handle both Western harmonic progressions and Indian melodic sensibilities.

Our method addresses this gap by incorporating the intuitive key-finding process used by Indian musicians, who rely on harmonic resonance rather than rigid theoretical frameworks.
This cultural grounding makes our approach particularly suitable for hybrid musical styles.

\section{Methodology}\label{sec:methodology}
Our approach consists of five main stages: audio preprocessing, harmonic overlay generation, hybrid signal creation, chroma feature extraction, and similarity-based evaluation.

\subsection{Audio Preprocessing}\label{subsec:audio-preprocessing}
We begin by loading the audio file and computing fundamental statistics.
The audio is normalized and converted to mono for ensuring consistent processing across different source materials.

\subsection{Harmonic Overlay Generation}\label{subsec:harmonic-overlay-generation}
For each of the 12 chromatic notes, we generate synthetic harmonic content.
The frequency of each candidate tonic is calculated using:
\begin{equation}
f_0 = 440 \cdot 2^{(n-69)/12}\label{eq:equation3}
\end{equation}
Where $n$ is the MIDI note number. We then generate a harmonic series:
\begin{equation}
s(t) = G \cdot \bar{A} \sum_{i=1}^{N} \frac{1}{i} \sin(2\pi \cdot i \cdot f_0 \cdot t)\label{eq:equation4}
\end{equation}
Where $G$ is the synthesis gain (0.8), $\bar{A}$ is the mean amplitude, and $N$ is the number of harmonics (6).


\subsection{Chroma Feature Extraction}\label{subsec:chroma-feature-extraction}
Chroma features represent the intensity of each pitch class (C, C\#, D, \ldots, B) independent of octave, making them ideal for harmonic analysis. We compute chroma features using the Constant-Q Transform:
\begin{equation}
C[k,n] = |CQT(h(t))[k,n]|^2\label{eq:equation2}
\end{equation}
Where $h(t)$ is the hybrid signal (original + synthetic), and chroma vectors are obtained by mapping CQT bins to pitch classes modulo 12.

\begin{equation}
\mathrm{Chroma}[p,n] =
\frac{\sum_{k \in \mathcal{K}(p)} \left| CQT[k,n] \right|^2}
{\sum_{p'=0}^{11} \sum_{k \in \mathcal{K}(p')} \left| CQT[k,n] \right|^2} \label{eq:chroma_cqt_norm}
\end{equation}

\subsection{Similarity-Based Evaluation}\label{subsec:similarity-based-evaluation}

To identify key, we employ a similarity-based strategy that compares chroma features of the original audio with those of the hybrid signals.
The audio is divided into fixed-length segments of $T_{\text{chunk}} = 20$ seconds, and for each segment a mean chroma vector is computed.
This chunked representation ensures temporal robustness by averaging local chroma variations while preserving tonal structure.

\paragraph{Ranking Candidates via Cosine Similarity:}
This similarity-based evaluation provides a principled way to rank candidate keys: the tonic whose hybrid signal achieves the highest mean cosine similarity with the original audio is selected as the most likely key.

\subsection{Iterative Refinement Based on Chroma Analysis}\label{subsec:iterative-refinement}
Initial experiments using tonic-only harmonic overlays yielded moderate success, with key estimation accuracy ranging between 60--70\%.
To investigate ranking inconsistencies, we analyzed \texttt{chroma\_cqt} plots and observed recurring tonal interference patterns.
In particular, strong harmonic energy at the interval of a perfect fifth often led to misclassification, as the fifth note would occasionally dominate the ranking.
To mitigate this, we extended the synthetic overlay to include both the candidate tonic and its perfect fifth.
This refinement was implemented by generating a secondary harmonic series for the fifth note and blending it with the tonic overlay at a reduced gain.
The resulting hybrid signal improved tonal coherence and consistently elevated the correct key to the top rank.

\subsection{Perfect Fifth Enhancement}\label{subsec:perfect-fifth-enhancement}
To address perfect fifth ambiguity, we simultaneously generate harmonics for both the candidate tonic and its perfect fifth:
\begin{equation}
s_{total}(t) = s_{tonic}(t) + \alpha \cdot s_{fifth}(t)\label{eq:equation}
\end{equation}
Where $\alpha = 0.75$ weights the perfect fifth contribution.

Analyzing the \texttt{chroma\_cqt} representation of this hybrid signal reveals that the chroma feature corresponding to the candidate tonic approaches unity, indicating strong tonal reinforcement. When the synthetic harmonic series for the perfect fifth is added, the resulting chroma vector also exhibits elevated energy at the fifth pitch class. This dual reinforcement ensures that the tonic remains dominant, while the fifth contributes supportive but non-overpowering energy.

The key insight emerges when cosine similarity is computed between the original audio and each hybrid signal. For the correct candidate tonic, the inclusion of its perfect fifth enhances tonal coherence, yielding a higher similarity score. Conversely, if the perfect fifth itself were incorrectly assumed as the candidate tonic, its own fifth would correspond to a different pitch class, thereby reducing the cosine similarity. This ranking effect systematically elevates the true tonic above its fifth, resolving the ambiguity observed in earlier tonic-only overlays.

\section{Empirical Analysis}\label{sec:empirical-analysis}

We present detailed analysis of a few music tracks to demonstrate the effectiveness of our perfect fifth enhancement approach.
Table \ref{tab:results-summary} summarizes the key findings across all test tracks.

\begin{table}[H]
\centering
\renewcommand{\arraystretch}{1.2}
\begin{tabular}{@{}lllll@{}}
\toprule
\textbf{Track} & \textbf{KS} & \textbf{FFT} & \textbf{Without PF} & \textbf{Ground Truth}  \\
\midrule
Music A & G Minor & C3 & F & F \\
\rowcolor{lightgray} Music B & E Minor & G\#2 & G\#  & G\# \\
\rowcolor{blue!10} Music C & G\# Major & F\#1 & B  & E \\
Music D & F\# Minor & F\#4 & A & A \\
Music E & D Major & F2 & A\# & A\# \\
Music F & B Minor & D\#1 & A & A \\
\rowcolor{blue!10} Music G & C Major & C4 & F & F \\
\rowcolor{blue!10} Music H & D Major & A\# & F & A\# \\
\rowcolor{lightgray} Music I & C Major & G1 & G & G \\
\rowcolor{lightgray} Music J & B Major & C\#4 & C\# & C\# \\
Music K & C\# Major & E2 & B & B \\
Music L & E Major & B1 & G\# & G\# \\
Music M & B Major & F2 & C\# & C\# \\
\bottomrule
\end{tabular}
\caption{Track-wise Summary of Key Detection Results}
\label{tab:results-summary}
\end{table}

The following observations summarize the performance of each key detection method across the evaluated tracks:

\begin{itemize}
    \item \textbf{Krumhansl-Schmuckler (KS) algorithm:} This approach shows frequent mismatches, particularly in tracks exhibiting modal ambiguity or non-functional harmony.
    Its reliance on Western tonal templates limits its effectiveness in diverse musical contexts.

    \item \textbf{FFT-Based Method:} The FFT-based approach estimates the dominant frequency by computing the Fast Fourier Transform (FFT) of the vocal signal and identifying the highest spectral
    Still, it gave better results than Krumhansl-Schmuckler algorithm, as highlighted in lightgray rows.

    \item \textbf{Synthetic Overlay without Perfect Fifth:} The overlay-only method performs well in simpler harmonic contexts but struggles when tonal ambiguity is present.
    Its performance degrades in compositions with competing harmonic centers or strong intervallic reinforcement.
\end{itemize}

\subsection{Music B}\label{subsec:song:-apni-to-jaise-taise}
\begin{figure}[H]
\centering
\safeincludegraphics[width=\columnwidth]{track_b_original.png}
\caption{Chromagram of the original Music B.}
\label{fig:apni-original}
\end{figure}

The image in Figure~\ref{fig:apni-original} shows the chromagram of the original track Music B.
You can see that the chroma corresponding to G\# has the highest intensity throughout the track, indicating that G\# is likely the tonic.

\begin{figure}[H]
\centering
\safeincludegraphics[width=\columnwidth]{no_perfect_fifth/track_b_G_sharp_hybrid.png}
\caption{Chromagram of the Music B with no perfect fifth overlay.}
\label{fig:apni-no-pf}
\end{figure}
\begin{table}[H]
\centering
\begin{tabular}{@{}lc@{}}
\toprule
\textbf{Note} & \textbf{Score} \\
\midrule
G\# & 0.9578 \\
B & 0.9091 \\
A & 0.9030 \\
\bottomrule
\end{tabular}
\caption{Cosine Similarity Scores for Leading Candidate Notes in without perfect fifth overlay ``Music B''.}
\label{tab:Track-B-No-PF-results}
\end{table}

As you can see in Table~\ref{tab:Track-B-No-PF-results}, G\# has the highest cosine similarity score of 0.9578, confirming it as the top candidate for the tonic.

\begin{figure}[H]
\centering
\safeincludegraphics[width=\columnwidth]{perfect_fifth/track_b_G_sharp_hybrid.png}
\caption{Chromagram of the Music B with perfect fifth overlay.}
\label{fig:apni-pf}
\end{figure}
\begin{table}[H]
\centering
\begin{tabular}{@{}lc@{}}
\toprule
\textbf{Note} & \textbf{Score} \\
\midrule
G\# & 0.9208 \\
B & 0.8806 \\
C\# & 0.8804 \\
\bottomrule
\end{tabular}
\caption{Cosine Similarity Scores for Leading Candidate Notes in with perfect fifth overlay ``Music B''.}
\label{tab:Track-B-PF-results}
\end{table}

This example illustrates that the key identification using perfect fifth and without perfect fifth enhancement yields the same top key, G\#.

\subsection{Music C}\label{subsec:song:-durge-durgat-bhari}

\begin{figure}[H]
\centering
\safeincludegraphics[width=\columnwidth]{track_c_original.png}
\caption{Chromagram of the original Music C.}
\label{fig:durge-original}
\end{figure}
The chromagram shows consistently high intensity in the E pitch class across time, suggesting that E is the tonal center or tonic of the piece.

\begin{table}[H]
\centering
\begin{tabular}{@{}lc@{}}
\toprule
\textbf{Note} & \textbf{Score} \\
\midrule
B & 0.8649 \\
E & 0.8636 \\
G\# & 0.8340 \\
\bottomrule
\end{tabular}
\caption{Cosine Similarity Scores for Leading Candidate Notes in without the perfect fifth overlay ``Music C''.}
\label{tab:Track-C-No-PF-results}
\end{table}

This demonstrates a classic case of perfect fifth confusion, as shown in Table~\ref{tab:Track-C-No-PF-results}.
Without the perfect fifth enhancement, the algorithm incorrectly identifies B as the tonic, while the actual key E ranks second with a marginally lower score.
This near-tie occurs because B is the perfect fifth of E, and the track likely contains E major or minor chords that emphasize both pitches.

\begin{table}[H]
\centering
\begin{tabular}{@{}lc@{}}
\toprule
\textbf{Note} & \textbf{Score} \\
\midrule
E & 0.8805 \\
B & 0.8516 \\
A & 0.8441 \\
\bottomrule
\end{tabular}
\caption{Cosine Similarity Scores for Leading Candidate Notes in with the perfect fifth overlay ``Music C''.}
\label{tab:Track-C-PF-results}
\end{table}


When perfect fifth harmonics are included, the true key E achieves the highest score (0.8805), creating a clearer separation from the confused fifth.

\subsection{Music G}\label{subsec:song:-hum-tumhe-chahate-hai}

\begin{figure}[H]
\centering
\safeincludegraphics[width=\columnwidth]{track_g_original.png}
\caption{Chromagram of the original Music G.}
\label{fig:hum-original}
\end{figure}

This music track has a long monologue section at the beginning with minimal harmonic content, making key identification challenging.
This monologue is more than 1 minute long, as shown in Figure~\ref{fig:hum-original}.

\begin{table}[H]
\centering
\begin{tabular}{@{}lc@{}}
\toprule
\textbf{Note} & \textbf{Score} \\
\midrule
G & 0.8408 \\
F & 0.8292 \\
A\# & 0.8176 \\
\bottomrule
\end{tabular}
\caption{Cosine Similarity Scores for Leading Candidate Notes in without the perfect fifth overlay ``Music G''.}
\label{tab:Track-G-No-PF-results}
\end{table}

This is another example of a perfect fifth ambiguity, as detailed in the Table~\ref{tab:Track-G-No-PF-results}.
The actual key F initially ranks second behind G.

\begin{table}[H]
\centering
\begin{tabular}{@{}lc@{}}
\toprule
\textbf{Note} & \textbf{Score} \\
\midrule
F & 0.8393 \\
C & 0.8342 \\
A\# & 0.8190 \\
\bottomrule
\end{tabular}
\caption{Cosine Similarity Scores for Leading Candidate Notes in with the perfect fifth overlay ``Music G''.}
\label{tab:Track-G-PF-results}
\end{table}

The perfect fifth enhancement successfully promotes F to first place (0.8393), while C (the perfect fifth) appears in second place (0.8343).
This ranking correctly reflects the tonic-dominant relationship in F major/minor tonality.

\subsection{Music H}\label{subsec:song:-kaise-muze-tum-mil-gaye}

\begin{figure}[H]
\centering
\safeincludegraphics[width=\columnwidth]{track_h_original.png}
\caption{Chromagram of the original track Music H.}
\label{fig:track-h-original}
\end{figure}

A\# shows high, recurrent intensity consistent with a tonic center.
F is nearly as strong, which is the perfect fifth of A\#.

\begin{table}[H]
\centering
\begin{tabular}{@{}lc@{}}
\toprule
\textbf{Note} & \textbf{Score} \\
\midrule
F & 0.9082 \\
G & 0.9062 \\
A\# & 0.9009 \\
\bottomrule
\end{tabular}
\caption{Cosine Similarity Scores for Leading Candidate Notes in without the perfect fifth overlay ``Music H''.}
\label{tab:Track-H-No-PF-results}
\end{table}

In this case, as shown in Table~\ref{tab:Track-H-No-PF-results}, when perfect fifth harmonics are not included, the actual key A\# ranks third with a score of 0.9009.
The perfect fifth of A\# is F, which initially achieves the highest score.
This demonstrates how the perfect fifth confusion can completely misidentify the tonic.

\begin{table}[H]
\centering
\begin{tabular}{@{}lc@{}}
\toprule
\textbf{Note} & \textbf{Score} \\
\midrule
A\# & 0.9100 \\
G & 0.9021 \\
F & 0.9015 \\
\bottomrule
\end{tabular}
\caption{Cosine Similarity Scores for Leading Candidate Notes in with the perfect fifth overlay ``Music G''.}
\label{tab:Track-H-PF-results}
\end{table}


With the perfect fifth enhancement, A\# correctly achieves first place (0.9100), while its perfect fifth F drops to third place (0.9015).
This reordering demonstrates the method's effectiveness in resolving harmonic ambiguities.

\section{Discussion}\label{sec:discussion}

\subsection{Theoretical Foundation}\label{subsec:theoretical-foundation}
Our approach is grounded in the acoustic principle that harmonically related tones will exhibit constructive interference when correctly aligned.
By adding synthetic harmonics at the true tonic frequency, we amplify existing harmonic content in the audio, creating stronger chroma representations.
The refinement emerged from a data-driven observation rather than a theoretical prescription.
By visualizing chroma\_cqt plots and analyzing ranking inconsistencies, we identified the perfect fifth as a recurring source of confusion.
This insight aligns with known harmonic reinforcement patterns in tonal music and demonstrates the value of exploratory analysis in MIR algorithm design.

\subsection{Limitations and Future Work}\label{subsec:limitations-and-future-work}
While the proposed method demonstrates promising results, several limitations remain:

\begin{itemize}
    \item \textbf{Single-key assumption:} The algorithm assumes a stable key throughout the entire track and does not handle key modulations or polytonality.

    \item \textbf{Genre-specific evaluation:} The current evaluation is limited to a narrow set of musical styles.
    Broader testing across diverse genres is needed to assess generalizability.

    \item \textbf{Chromatic and atonal complexity:} Highly chromatic or atonal compositions may exhibit ambiguous tonal centers, reducing the effectiveness of harmonic overlays.

    \item \textbf{Fixed parameter settings:} Parameters such as synthesis gain, number of harmonics, and weighting coefficients are fixed across all inputs. Adaptive tuning based on genre or signal characteristics could improve performance.
\end{itemize}

Future work will explore dynamic parameter optimization, support for multi-key detection, and genre-aware harmonic modeling to enhance robustness and scalability.

Future research directions include adaptive parameter tuning based on musical style classification, extension to modulation detection and tracking, integration with advanced pitch tracking systems like CREPE, expansion to other world music traditions, and real-time streaming implementation.

\section{Conclusion}\label{sec:conclusion}
This paper presents a unique approach to key identification in music that combines synthetic harmonic overlays with chroma analysis.
Our method successfully addresses the challenge of perfect fifth ambiguity, a common problem in chord-rich musical genres.

Key contributions include: (1) a culturally adaptive algorithm that emulates Indian musicians' intuitive key-finding process, (2) a solution to perfect fifth confusion using a dual harmonic overlay, (3) empirical validation showing consistent improvements in key identification accuracy, and (4) computational efficiency that is suitable for practical applications.

The method achieves an average accuracy improvement of 1.3\% with significantly enhanced confidence scores.
This work bridges cultural musical practices with computational analysis, offering a blueprint for developing culturally sensitive MIR systems in the future.

\balance

\section*{Acknowledgments}
Portions of this manuscript were drafted or refined using generative AI tools (Microsoft CoPilot).
All content was reviewed and edited by the authors to ensure accuracy, originality, and scholarly integrity.
The authors take full responsibility for the final version of the manuscript.

\begin{thebibliography}{99}
\bibitem{krumhansl1990cognitive}
C. L. Krumhansl, \textit{Cognitive Foundations of Musical Pitch}. Oxford University Press, 1990.

\bibitem{kim2018crepe}
J. W. Kim, J. Salamon, P. Li, and J. P. Bello, ``CREPE: A convolutional representation for pitch estimation,'' in \textit{Proc. ICASSP}, 2018, pp. 161-165.

\bibitem{muller2011chroma}
M. M\"uller and S. Ewert, ``Chroma toolbox: MATLAB implementations for extracting variants of chroma-based audio features,'' in \textit{Proc. ISMIR}, 2011, pp. 215-220.

\bibitem{srinivasamurthy2021saraga}
A. Srinivasamurthy, S. Gulati, R. C. Repetto, and X. Serra, ``Saraga: Open datasets for research on Indian art music,'' \textit{Empirical Musicology Review}, vol. 16, no. 1, pp. 85-98, 2021.

\bibitem{raja2023sequential}
V. N. Senthil Raja, ``Sequential pitch distributions for raga detection,'' in \textit{Proc. AI Music Creativity Conference}, 2023.

\bibitem{silva2025}
Henrique Perez Gomes da Silva, \textit{Understanding the Algorithm Behind Audio Key Detection}, arXiv preprint\href{https://arxiv.org/abs/2505.17259}{arXiv:2505.17259}, 2025.

\end{thebibliography}

\end{document}
